% Luc Destombe --- licence CC BY-NC-SA 4.0
% https://creativecommons.org/licenses/by-nc-sa/4.0/deed.fr
% Fichier autonome : compiler avec pdflatex, deux fois (signets, nombre de pages).
\def\lyceeversion{1.0}
\documentclass[a4paper,10pt]{article}
\makeatletter
\newif\iflycee@niveaux
\lycee@niveauxtrue
\newif\iflycee@palatino
\lycee@palatinotrue

\RequirePackage[utf8]{inputenc}
\RequirePackage[T1]{fontenc}
\RequirePackage[french]{babel}
\RequirePackage{etoolbox}

\newcommand{\ShorthandsOff}{\@ifundefined{shorthandoff}{}{\shorthandoff{;:!?}}}
\newcommand{\ShorthandsOn}{\@ifundefined{shorthandon}{}{\shorthandon{;:!?}}}
\ShorthandsOff
\AtBeginDocument{\ShorthandsOn}
\AtBeginEnvironment{tikzpicture}{\ShorthandsOff}

\RequirePackage{geometry}
\RequirePackage{fancyhdr}
\RequirePackage{lastpage}
\RequirePackage{multicol}
\RequirePackage{array}
\RequirePackage{enumitem}
\RequirePackage{xcolor}
\RequirePackage{needspace}

\RequirePackage{amsmath,amssymb}
\RequirePackage{mathpazo}
\renewcommand{\ttdefault}{lmtt}
\DeclareMathSizes{10.95}{10.95}{8}{5.5}
\begingroup\catcode`\_=\active \catcode`\^=\active
\gdef_#1{\sb{\scriptscriptstyle #1}}
\gdef^#1{\sp{\scriptscriptstyle\lycee@moinsepais #1}}
\endgroup
\newcommand\lycee@moins{\mathbin{%
  \setbox\z@\hbox{$\m@th\scriptscriptstyle\lycee@moinsorig$}%
  \dimen@=.55\wd\z@ \advance\dimen@-.5pt
  \hbox to.75\wd\z@{\hss$\m@th\scriptscriptstyle\vcenter{\hbox{%
    \tikz\draw[line width=.5pt,line cap=round](0,0)--(\the\dimen@,0);}}$\hss}}}
\begingroup\lccode`\~=`\- \lowercase{\endgroup\let~}\lycee@moins
\newcommand\lycee@moinsepais{\mathcode`\-="8000 }
\AtBeginDocument{\mathchardef\lycee@moinsorig=\mathcode`\-\relax}
\AtBeginDocument{\catcode`\_=12 \mathcode`\_="8000
  \catcode`\^=12 \mathcode`\^="8000 }
\newcommand\lycee@exposants{%
  \fontdimen13\textfont2=.48\fontdimen6\textfont2
  \fontdimen14\textfont2=.48\fontdimen6\textfont2
  \fontdimen15\textfont2=.48\fontdimen6\textfont2
  \fontdimen13\scriptfont2=.48\fontdimen6\scriptfont2
  \fontdimen14\scriptfont2=.48\fontdimen6\scriptfont2
  \fontdimen15\scriptfont2=.48\fontdimen6\scriptfont2}
\every@math@size\expandafter{\the\every@math@size\lycee@exposants}
\newlength{\retraitcalculs}
\setlength{\retraitcalculs}{2em}

\RequirePackage{tikz}
\usetikzlibrary{arrows.meta,calc,babel,shapes.misc}
\RequirePackage[most]{tcolorbox}
\tcbuselibrary{skins,breakable}

\definecolor{coulDef}{HTML}{1F4E79}
\definecolor{coulProp}{HTML}{7030A0}
\definecolor{coulRem}{HTML}{C55A11}
\definecolor{coulEx}{HTML}{2E7D32}
\definecolor{coulExo}{HTML}{B03A2E}
\definecolor{coulPart}{HTML}{1F4E79}
\definecolor{coulMeth}{HTML}{1B6E4F}
\definecolor{coulCourbe}{HTML}{0B5ED7}
\definecolor{coulCourbeB}{HTML}{C00000}
\definecolor{coulCourbeC}{HTML}{2E7D32}
\definecolor{coulGrille}{HTML}{8C8C8C}
\definecolor{coulRep}{HTML}{1B6E4F}
\definecolor{coulBar}{HTML}{2E7D32}

\newcommand{\R}{\mathbb{R}}
\newcommand{\Rs}{\mathbb{R}^{*}}
\renewcommand{\le}{\leqslant}
\renewcommand{\ge}{\geqslant}
\renewcommand{\approx}{\simeq}
\newcommand{\vect}[1]{\overrightarrow{#1}}
\newcommand{\norme}[1]{\left\lVert #1 \right\rVert}
\newcommand{\intervalle}[2]{\left[#1\,;#2\right]}
\RequirePackage{cancel}
\renewcommand{\CancelColor}{\color{coulExo}}
\newcommand{\fois}[1]{\times{\color{coulExo}#1}}
\newcommand{\barre}[1]{\cancel{#1}}
\newlength{\lycee@eqh}\newlength{\lycee@eqd}
\newsavebox{\lycee@eqa}\newsavebox{\lycee@eqb}\newsavebox{\lycee@eqc}
\newcommand{\lycee@eqcol}[1]{\hbox{$\renewcommand{\arraystretch}{1.35}%
  \begin{array}[t]{@{}>{\displaystyle}r@{}>{\displaystyle{}}l@{}}#1\end{array}$}}
\newcommand{\lycee@eqmesure}[1]{%
  \ifdim\ht#1>\lycee@eqh\setlength{\lycee@eqh}{\ht#1}\fi
  \ifdim\dp#1>\lycee@eqd\setlength{\lycee@eqd}{\dp#1}\fi}
\newcommand{\lycee@eqtrait}{\hspace{1.2em}\tikz[baseline=0pt]\draw[dash pattern=on 2pt off 2pt,
  line width=0.5pt] (0,\lycee@eqh) -- (0,-\lycee@eqd);\hspace{1.2em}}
\newcommand{\eqproduit}[2]{\par\smallskip\noindent C'est une équation produit :
  $A\times B=0 \iff A=0$ ou $B=0$.\par\nopagebreak\smallskip
  \begingroup\setlength{\lycee@eqh}{0pt}\setlength{\lycee@eqd}{0pt}%
  \setbox\lycee@eqa=\lycee@eqcol{#1}\setbox\lycee@eqb=\lycee@eqcol{#2}%
  \lycee@eqmesure\lycee@eqa\lycee@eqmesure\lycee@eqb%
  \noindent\hspace*{\retraitcalculs}\usebox\lycee@eqa\lycee@eqtrait\usebox\lycee@eqb\endgroup\par\smallskip}
\newcommand{\eqproduittrois}[3]{\par\smallskip\noindent C'est une équation produit :
  $A\times B\times C=0 \iff A=0$ ou $B=0$ ou $C=0$.\par\nopagebreak\smallskip
  \begingroup\setlength{\lycee@eqh}{0pt}\setlength{\lycee@eqd}{0pt}%
  \setbox\lycee@eqa=\lycee@eqcol{#1}\setbox\lycee@eqb=\lycee@eqcol{#2}\setbox\lycee@eqc=\lycee@eqcol{#3}%
  \lycee@eqmesure\lycee@eqa\lycee@eqmesure\lycee@eqb\lycee@eqmesure\lycee@eqc%
  \noindent\hspace*{\retraitcalculs}\usebox\lycee@eqa\lycee@eqtrait\usebox\lycee@eqb\lycee@eqtrait\usebox\lycee@eqc\endgroup
  \par\smallskip}

\tikzset{grille/.style={coulGrille,line width=0.4pt}}
\newcommand{\quadrillage}[4]{\draw[grille,step=1] (#1,#2) grid (#3,#4);}
\tikzset{
  croix/.style={cross out,draw,line width=0.6pt,inner sep=0pt,minimum size=4pt},
  croixdroite/.style={croix,rotate=45,line width=0.8pt},
}
\newcommand{\pt}[3]{\path (#1) node[croix]{} node[#2,font=\small,inner sep=2.5pt]{$#3$};}

\newcommand{\lycee@repere}[4]{%
  \draw[grille,step=1] (#1,#2) grid (#3,#4);
  \draw[-{Stealth[length=2mm]},thick] (#1,0)--({#3+0.4},0) node[below left=-1pt]{$x$};
  \draw[-{Stealth[length=2mm]},thick] (0,#2)--(0,{#4+0.4}) node[below left=-1pt]{$y$};
  \node[below left,font=\scriptsize,inner sep=1.5pt] at (0,0) {$0$};}
\newcommand{\lycee@gradx}[1]{\foreach \k/\l in {#1}{%
  \draw (\k,0.07)--(\k,-0.07) node[below,font=\scriptsize,inner sep=1.5pt]{$\l$};}}
\newcommand{\lycee@grady}[1]{\foreach \k/\l in {#1}{%
  \draw (0.07,\k)--(-0.07,\k) node[left,font=\scriptsize,inner sep=1.5pt]{$\l$};}}
\newcommand{\lycee@intff}[2]{\left[#1\,;#2\right]}
\newcommand{\lycee@intfo}[2]{\left[#1\,;#2\right[}
\newcommand{\lycee@intof}[2]{\left]#1\,;#2\right]}
\newcommand{\lycee@intoo}[2]{\left]#1\,;#2\right[}
\newcommand{\lycee@ens}[1]{\left\{#1\right\}}
\AtBeginDocument{%
  \providecommand{\repere}{\lycee@repere}%
  \providecommand{\gradx}{\lycee@gradx}%
  \providecommand{\grady}{\lycee@grady}%
  \providecommand{\Cf}{\mathcal{C}\sb{\scriptscriptstyle f}}%
  \providecommand{\Cg}{\mathcal{C}\sb{\scriptscriptstyle g}}%
  \providecommand{\intff}{\lycee@intff}%
  \providecommand{\intfo}{\lycee@intfo}%
  \providecommand{\intof}{\lycee@intof}%
  \providecommand{\intoo}{\lycee@intoo}%
  \providecommand{\ens}{\lycee@ens}}

\newlist{questions}{enumerate}{3}
\setlist[questions,1]{label=\arabic*\textdegree),leftmargin=*,itemsep=2pt,topsep=3pt}
\setlist[questions,2]{label=\alph*),leftmargin=*,itemsep=1pt,topsep=2pt}
\setlist[questions,3]{label=\roman*),leftmargin=*,itemsep=1pt,topsep=1pt}
\setlist[itemize]{label=\textbullet,leftmargin=*,itemsep=2pt,topsep=3pt}

\newcommand{\besoinplace}[1]{\par\penalty\@M\begingroup
  \dimen@=#1\relax\dimen@ii=\pagegoal\advance\dimen@ii-\pagetotal
  \ifdim\dimen@>\dimen@ii\ifdim\dimen@ii>\z@\vfil\fi\break\fi\endgroup}

\newcommand{\lycee@pied}{\thepage/\pageref{LastPage}}
\newcommand{\entete}[2]{%
  \pagestyle{fancy}\fancyhf{}%
  \fancyhead[L]{\footnotesize\color{coulPart}\textbf{#1}}%
  \fancyhead[R]{\footnotesize\color{coulPart}\textbf{#2}}%
  \fancyfoot[C]{\footnotesize\lycee@pied}%
  \renewcommand{\headrulewidth}{0.4pt}%
  \renewcommand{\headrule}{\hbox to\headwidth{%
    \color{coulPart!40}\leaders\hrule height \headrulewidth\hfill}}}

\RequirePackage{ccicons}
\newcommand{\lycee@licence}{{\scriptsize\color{gray}%
  \copyright\ Luc Destombe --- \ccbyncsa\ CC BY-NC-SA 4.0}}
\newif\iflycee@licence \lycee@licencetrue
\newcommand{\sanslicence}{\lycee@licencefalse}
\AtBeginDocument{\iflycee@licence\AddToHookNext{shipout/foreground}{%
  \put(0,0){\rlap{\kern\dimexpr1in+\hoffset+\oddsidemargin+\textwidth\relax
    \llap{\raisebox{-\dimexpr1in+\voffset+\topmargin+\headheight+\headsep
      +\textheight+\footskip\relax}{\lycee@licence}}}}}\fi}

\newcommand{\lycee@site}{faire-des-maths.fr}
\newcommand{\lycee@version}{{\scriptsize\color{gray}\lycee@site
  \ifdefined\lyceeversion\ --- Version~\lyceeversion\fi}}
\AtBeginDocument{\iflycee@licence\AddToHookNext{shipout/foreground}{%
  \put(0,0){\rlap{\kern\dimexpr1in+\hoffset+\oddsidemargin\relax
    \raisebox{-\dimexpr1in+\voffset+\topmargin+\headheight+\headsep
      +\textheight+\footskip\relax}{\lycee@version}}}}\fi}

\newcommand{\formulesserrees}{%
  \setlength{\abovedisplayskip}{3pt}\setlength{\belowdisplayskip}{3pt}%
  \setlength{\abovedisplayshortskip}{2pt}\setlength{\belowdisplayshortskip}{3pt}}

\setlength{\parindent}{0pt}

\RequirePackage[implicit=false,hidelinks,pdfencoding=auto,psdextra,bookmarksopen,
  bookmarksopenlevel=1,pdfcreator={},pdfproducer={}]{hyperref}
\RequirePackage{bookmark}
\hypersetup{pdfauthor={Luc Destombe},pdfkeywords={CC BY-NC-SA 4.0}}
\newcounter{lycee@signet}
\newcommand{\signet}[3][]{\stepcounter{lycee@signet}%
  \edef\lycee@dest{\if\relax\detokenize{#1}\relax
    signet.\arabic{lycee@signet}\else#1\fi}%
  \hypertarget{\lycee@dest}{}\bookmark[level=#2,dest=\lycee@dest]{#3}}
\pdfstringdefDisableCommands{%
  \def\R{R}\def\vect#1{#1}\def\Cf{Cf}\def\Cg{Cg}\def\fois#1{×#1}%
  \def\barre#1{#1}\def\intervalle#1#2{[#1;#2]}\def\intff#1#2{[#1;#2]}%
  \def\textsuperscript#1{#1}\def\dfrac#1#2{#1/#2}\def\frac#1#2{#1/#2}%
  \def\vec#1{#1⃗}}%
\RequirePackage[official]{eurosym}

\geometry{top=0.8cm,bottom=1.3cm,left=1cm,right=1cm,footskip=0.6cm}

\pagestyle{fancy}
\fancyhf{}
\renewcommand{\headrulewidth}{0pt}
\fancyfoot[C]{\footnotesize Page \thepage{} / \pageref{LastPage}}

\newcommand{\titrefiche}[2]{%
  \begin{center}
    \Large\textbf{#1}\\[2pt]
    \large\textbf{#2}\\[2pt]
  \end{center}
  \vspace{0.10cm}}

\setlength{\columnseprule}{0.4pt}
\setlength{\columnsep}{15pt}
\renewcommand{\columnseprulecolor}{\color{black!45}}
\setlength{\multicolsep}{2pt}

\newcounter{partiefiche}
\renewcommand{\thepartiefiche}{\Roman{partiefiche}}
\newif\ifapresbandeau
\newcommand{\lycee@bandeau}[1]{%
  \par\penalty-600
  \vspace{0.08cm}%
  \noindent\colorbox{coulPart}{%
    \parbox{\dimexpr\linewidth-2\fboxsep\relax}{%
      \color{white}\bfseries\raggedright\strut#1\strut}}%
  \nopagebreak\par\nopagebreak
  \vspace{0.06cm}\nopagebreak
  \apresbandeautrue}
\newcommand{\partiefiche}[1]{%
  \refstepcounter{partiefiche}%
  \lycee@bandeau{\signet{1}{\thepartiefiche{} --- #1}\thepartiefiche{} --- #1}}
\newcommand{\partiefichesansnum}[1]{\lycee@bandeau{\signet{1}{#1}#1}}

\newcounter{exo}
\newlength{\espaceexo}\setlength{\espaceexo}{7pt}
\newlength{\espacetitre}\setlength{\espacetitre}{2pt}
\newcommand{\exo}[1][\relax]{%
  \par
  \ifapresbandeau\apresbandeaufalse\else\penalty-150\addvspace{\espaceexo}\fi
  \refstepcounter{exo}%
  \noindent\signet[exercice-\arabic{exo}]{2}{Exercice \arabic{exo}}%
  \textbf{\color{coulExo}\underline{Exercice \theexo}}%
  \ifx\relax#1\relax\else\textbf{ : #1}\fi
  \par\nopagebreak\vspace{\espacetitre}\nopagebreak
  \ignorespaces}

\setlist[enumerate]{nosep,leftmargin=*,itemsep=2pt,topsep=2pt}
\setlist[itemize]{nosep,leftmargin=*,topsep=2pt}
\newcommand{\choix}[3]{\par\hspace*{1em}%
  \makebox[0.3\linewidth][l]{a) #1}\makebox[0.3\linewidth][l]{b) #2}c) #3}
\newcommand{\deux}[2]{\par\hspace*{0.6em}%
  \makebox[0.48\linewidth][l]{#1}#2}
\newcommand{\trois}[3]{\par\hspace*{0.6em}%
  \makebox[0.32\linewidth][l]{#1}\makebox[0.32\linewidth][l]{#2}#3}

  \renewcommand{\exo}[1][\relax]{%
    \ifapresbandeau\apresbandeaufalse\else\par\penalty-150\fi
    \refstepcounter{exo}%
    \vspace{0.05cm}%
    \noindent\signet[exercice-\arabic{exo}]{2}{Exercice \arabic{exo}}%
    \textbf{\color{coulExo}\underline{Exercice \theexo}}%
    \ifx\relax#1\relax\else\textbf{ : #1}\fi%
    \par\nopagebreak
    \ignorespaces}
  \newcommand{\rappel}[1]{%
    \par\vspace{0.06cm}\nopagebreak
    \noindent\fcolorbox{black!35}{black!5}{%
      \parbox{\dimexpr\linewidth-2\fboxsep-2\fboxrule\relax}{%
        \small\textbf{Rappel.} #1}}%
    \par\vspace{0.06cm}\nopagebreak}
  \setlist[enumerate]{nosep,leftmargin=*}
  \setlist[itemize]{nosep,leftmargin=*}
  \setlength{\multicolsep}{3pt plus 1pt minus 1pt}
  \setlength{\premulticols}{10pt}
  \setlength{\postmulticols}{10pt}
  \newlist{expr}{enumerate}{1}
  \setlist[expr]{label={\Alph*\ $=$},nosep,leftmargin=*,itemsep=0pt}

\renewenvironment{center}{\par\addvspace{3pt}\centering}{\par\addvspace{3pt}}
\formulesserrees
\AtBeginDocument{\formulesserrees}
\clubpenalty=10000
\widowpenalty=10000
\displaywidowpenalty=10000
\brokenpenalty=10000
\setlength{\parskip}{0pt}

\RequirePackage{ragged2e}
\setlength{\RaggedRightRightskip}{0pt plus 1fil}
\AtBeginDocument{\RaggedRight\binoppenalty=10000 \relpenalty=10000 }
\g@addto@macro\@parboxrestore{\RaggedRight}
\makeatother
\usepackage{tkz-tab}

\ShorthandsOff
\tikzset{lecture/.style={dashed,black!60,line width=0.6pt}}
\newcommand{\cerclenu}{%
  \draw[-{Stealth[length=2mm]},thick] (-1.3,0)--(1.3,0) node[below,font=\scriptsize,inner sep=2pt]{$x$};
  \draw[-{Stealth[length=2mm]},thick] (0,-1.3)--(0,1.3) node[right,font=\scriptsize,inner sep=2pt]{$y$};
  \draw[line width=0.9pt] (0,0) circle (1);
  \node[below left,font=\scriptsize,inner sep=1.5pt] at (0,0) {$O$};}
\newcommand{\cerclebase}{\cerclenu
  \node[below right,font=\scriptsize,inner sep=1.5pt] at (1,0) {$1$};
  \node[below left,font=\scriptsize,inner sep=1.5pt] at (-1,0) {$-1$};
  \node[above right,font=\scriptsize,inner sep=1.5pt] at (0,1) {$1$};
  \node[below right,font=\scriptsize,inner sep=1.5pt] at (0,-1) {$-1$};}
\newcommand{\cerclerep}{\cerclenu
  \node[above right,font=\scriptsize,inner sep=1.5pt] at (1,0) {$0$};
  \node[below left,font=\scriptsize,inner sep=1.5pt] at (1,0) {$2\pi$};
  \node[above left,font=\scriptsize,inner sep=1.5pt] at (-1,0) {$\pi$};}
\newcommand{\rayons}[1]{\foreach \a in {#1}{\draw[grille] (0,0)--({cos(\a)},{sin(\a)});}}
\newcommand{\ptc}[2]{%
  \path[coulCourbeB] ({cos(#1)},{sin(#1)}) node[croix]{};
  \node[font=\scriptsize,inner sep=1pt] at ({1.25*cos(#1)},{1.25*sin(#1)}) {#2};}
\newcommand{\reperetrigo}[2]{%
  \foreach \k in {#1,...,#2}{\draw[grille] ({\k*pi/6},-1.5)--({\k*pi/6},1.5);}
  \foreach \j in {-3,...,3}{\draw[grille] ({#1*pi/6},{\j/2})--({#2*pi/6},{\j/2});}
  \draw[-{Stealth[length=2mm]},thick] ({#1*pi/6},0)--({#2*pi/6+0.4},0)
    node[below left=-1pt,font=\footnotesize]{$x$};
  \draw[-{Stealth[length=2mm]},thick] (0,-1.5)--(0,1.8)
    node[below left=-1pt,font=\footnotesize]{$y$};
  \foreach \j/\l in {1/1,0.5/\frac{1}{2},-0.5/-\frac{1}{2},-1/-1}{%
    \node[left,font=\scriptsize,inner sep=1pt,outer sep=1pt,fill=white] at (0,\j) {$\l$};}}
\newcommand{\gradpi}[2][below]{\foreach \k/\l in {#2}{%
  \draw ({\k*pi/6},0.05)--({\k*pi/6},-0.05);
  \node[#1,font=\scriptsize,inner sep=1pt,outer sep=2pt,fill=white] at ({\k*pi/6},0) {$\l$};}}
\newcommand{\courbetrigo}[1]{\draw[coulCourbe,line width=1.1pt]
  plot[domain=-2*pi:2*pi,samples=121] (\x,{#1(\x r)});}
\ShorthandsOn

\begin{document}

\titrefiche{1\textsuperscript{re} spécialité --- Trigonométrie}{Fiche d'exercices du chapitre}

\begin{multicols*}{2}\raggedcolumns

\partiefiche{Cercle trigonométrique et radian}

\exo[Des degrés aux radians]
Convertir chaque angle en radians : a)~$40^{\circ}$ ; b)~$150^{\circ}$ ;
c)~$225^{\circ}$ ; d)~$300^{\circ}$ ; e)~$72^{\circ}$.

\exo[Des radians aux degrés]
Convertir chaque angle en degrés : a)~$\frac{3\pi}{4}$ ; b)~$\frac{7\pi}{12}$ ;
c)~$\frac{4\pi}{9}$ ; d)~$\frac{11\pi}{6}$ ; e)~$\frac{3\pi}{8}$.

\exo[Sans calculatrice]
Un angle de $15^{\circ}$ mesure $\frac{\pi}{12}$ rad. Sans calculatrice, en déduire la
mesure en radians d'un angle de $45^{\circ}$, de $75^{\circ}$, de $105^{\circ}$ et de
$165^{\circ}$.

\exo[Polygones réguliers]
Les sommets d'un polygone régulier sont sur un cercle de centre $O$ ; deux sommets
consécutifs forment avec $O$ un angle au centre.
\begin{enumerate}[label=\alph*)]
  \item Pour chaque polygone, triangle équilatéral, carré, hexagone, octogone et
        dodécagone ($12$ côtés), donner l'angle au centre en radians.
  \item L'angle au centre d'un polygone régulier mesure $24^{\circ}$. Combien a-t-il de
        côtés ? Donner cet angle en radians.
\end{enumerate}

\partiefiche{Longueur d'un arc, mesures d'un angle orienté}

\exo[Longueur d'un arc]
Sur un cercle de rayon $r$, un angle au centre de mesure $\alpha$ radians intercepte un
arc de longueur $\ell$. Calculer :
\begin{enumerate}[label=\alph*)]
  \item $\ell$, pour $r=6$ cm et $\alpha=\frac{\pi}{3}$ ;
  \item $\ell$, pour $r=9$ cm et $\alpha=\frac{5\pi}{6}$ ;
  \item $\alpha$, pour $r=10$ cm et $\ell=4$ cm ;
  \item $\alpha$, pour $r=4$ cm et $\ell=3\pi$ cm.
\end{enumerate}

\exo[Lire les mesures d'un point]
Pour chaque point $E$ et $F$ de la figure, lire sa mesure dans le sens direct, entre
$0$ et $2\pi$, puis sa mesure dans le sens indirect, entre $-2\pi$ et $0$.
\begin{center}
\begin{tikzpicture}[x=1cm,y=1cm]
  \rayons{0,45,...,315}\cerclerep\ptc{225}{$E$}
  \node[font=\footnotesize] at (0,1.65) {huitièmes de tour};
\end{tikzpicture}\hspace{0.6cm}
\begin{tikzpicture}[x=1cm,y=1cm]
  \rayons{0,30,...,330}\cerclerep\ptc{300}{$F$}
  \node[font=\footnotesize] at (0,1.65) {douzièmes de tour};
\end{tikzpicture}
\end{center}

\exo[Placer des images]
Sur le cercle ci-dessous (rayons tous les $\frac{\pi}{12}$), placer l'image de chaque
nombre réel :
\begin{center}
$A$ : $\frac{3\pi}{4}$ \quad $B$ : $-\frac{\pi}{6}$ \quad $C$ : $\frac{13\pi}{4}$ \quad
$D$ : $-\frac{7\pi}{3}$ \quad $G$ : $\frac{9\pi}{2}$ \quad $H$ : $-5\pi$
\end{center}
\begin{center}
\begin{tikzpicture}[x=1.3cm,y=1.3cm]
  \rayons{0,15,...,345}\cerclerep
\end{tikzpicture}
\end{center}

\exo[Même image ou non ?]
Pour chaque couple de nombres réels, dire s'ils ont la même image sur le cercle
trigonométrique, et rédiger la réponse par une phrase (« \ldots{} et \ldots{} ont la
même image, car \ldots{} »).
\begin{center}
a)~$7\pi$ et $-3\pi$ \quad b)~$\frac{3\pi}{4}$ et $\frac{19\pi}{4}$ \quad
c)~$\frac{5\pi}{2}$ et $-\frac{\pi}{2}$ \quad d)~$\frac{2\pi}{3}$ et $-\frac{16\pi}{3}$
\end{center}

\exo[Grande roue]
Une grande roue a un rayon de $30$ m. Une nacelle part du point le plus à droite de la
roue et tourne dans le sens direct.
\begin{enumerate}[label=\alph*)]
  \item La roue tourne de $\frac{5\pi}{2}$ rad. Quelle distance la nacelle
        parcourt-elle ? Donner la valeur exacte, puis une valeur approchée au mètre.
  \item Combien de tours la roue a-t-elle faits ? Où se trouve la nacelle ?
  \item Convertir en degrés un angle de $\frac{5\pi}{2}$ rad.
  \item La nacelle a parcouru $100\pi$ m. De quel angle la roue a-t-elle tourné ?
        De quel nombre réel, entre $0$ et $2\pi$, la nacelle est-elle alors l'image ?
\end{enumerate}

\partiefiche{Mesure principale d'un angle orienté}

\exo[Mesure principale]
Déterminer la mesure principale de chaque angle orienté, de mesure :
a)~$\frac{9\pi}{4}$ ; b)~$-\frac{7\pi}{3}$ ; c)~$\frac{23\pi}{6}$ ;
d)~$\frac{31\pi}{4}$ ; e)~$-\frac{21\pi}{2}$ ; f)~$\frac{199\pi}{3}$ ; g)~$18\pi$.

\exo[L'intrus]
Dans chaque liste, toutes les mesures sauf une sont celles d'un même angle orienté.
Trouver l'intrus, puis la mesure principale de l'angle.
\begin{enumerate}[label=\alph*)]
  \item $\frac{5\pi}{3}$ \quad $-\frac{\pi}{3}$ \quad $\frac{11\pi}{3}$ \quad
        $-\frac{7\pi}{3}$ \quad $\frac{7\pi}{3}$
  \item $\frac{3\pi}{4}$ \quad $-\frac{5\pi}{4}$ \quad $\frac{11\pi}{4}$ \quad
        $\frac{19\pi}{4}$ \quad $-\frac{3\pi}{4}$
\end{enumerate}

\exo[Satellite]
Un satellite tourne autour de la Terre, dans le sens direct, sur une orbite circulaire
de rayon $6\,000$ km. Depuis son lancement, il a tourné de $\frac{47\pi}{3}$ rad.
\begin{enumerate}[label=\alph*)]
  \item Combien de tours complets a-t-il faits ?
  \item Déterminer la mesure principale de sa position.
  \item Quelle distance a-t-il parcourue ? Donner la valeur exacte, puis une valeur
        approchée au kilomètre.
\end{enumerate}

\exo[Toutes les mesures d'un intervalle]
\begin{enumerate}[label=\alph*)]
  \item Donner la mesure principale de l'angle orienté de mesure $\frac{7\pi}{6}$.
  \item Donner tous les nombres réels de l'intervalle $\intff{-4\pi}{4\pi}$ qui ont la
        même image que $\frac{7\pi}{6}$.
\end{enumerate}

\partiefiche{Cosinus et sinus d'un nombre réel}

\exo[Valeurs exactes]
Déterminer la valeur exacte de : a)~$\cos\left(\frac{13\pi}{6}\right)$ ;
b)~$\sin\left(\frac{9\pi}{4}\right)$ ; c)~$\cos(-5\pi)$ ;
d)~$\sin\left(-\frac{7\pi}{2}\right)$ ; e)~$\cos\left(\frac{25\pi}{3}\right)$ ;
f)~$\sin\left(-\frac{15\pi}{4}\right)$.

\exo[Calculer une expression]
Calculer la valeur exacte de :
\begin{enumerate}[label=\alph*)]
  \item $f(x)=4\cos(x)-6\sin(x)$, pour $x=\frac{\pi}{6}$ ;
  \item $g(x)=2\sin^{2}(x)$, pour $x=\frac{\pi}{3}$ ;
  \item $h(x)=\sin(2x)$, pour $x=\frac{\pi}{12}$ ;
  \item $k(x)=\cos(x)\times\sin(x)$, pour $x=\frac{\pi}{4}$ ;
  \item $m(x)=\cos^{2}(x)-\sin^{2}(x)$, pour $x=\frac{\pi}{6}$.
\end{enumerate}

\exo[Le sinus d'après le cosinus]
À l'aide de la relation $\cos^{2}(x)+\sin^{2}(x)=1$, calculer la valeur exacte :
\begin{enumerate}[label=\alph*)]
  \item de $\sin(x)$, sachant que $\cos(x)=\frac{1}{4}$ et $x\in\intff{0}{\pi}$ ;
  \item de $\cos(x)$, sachant que $\sin(x)=-\frac{2}{3}$ et $x\in\intff{-\frac{\pi}{2}}{0}$ ;
  \item de $\sin(x)$, sachant que $\cos(x)=-\frac{12}{13}$ et $x\in\intff{-\pi}{0}$.
  \item Existe-t-il un nombre réel $x$ tel que $\cos(x)=\frac{2}{3}$ et $\sin(x)=\frac{2}{3}$ ?
\end{enumerate}

\exo[Démontrer les valeurs en $\frac{\pi}{4}$]
\begin{minipage}[c]{0.62\linewidth}\raggedright
$M$ est l'image de $\frac{\pi}{4}$ et $H$ son projeté sur l'axe des abscisses.
\begin{enumerate}[label=\alph*)]
  \item Quels sont les signes de $\cos\left(\frac{\pi}{4}\right)$ et de
        $\sin\left(\frac{\pi}{4}\right)$ ?
  \item Quelle est la nature du triangle $OHM$ ? En déduire que
        $\cos\left(\frac{\pi}{4}\right)=\sin\left(\frac{\pi}{4}\right)$.
  \item À l'aide de $\cos^{2}(x)+\sin^{2}(x)=1$, calculer
        $\cos\left(\frac{\pi}{4}\right)$.
\end{enumerate}
\end{minipage}\hfill
\begin{minipage}[c]{0.34\linewidth}\centering
\begin{tikzpicture}[x=2cm,y=2cm]
  \draw[-{Stealth[length=2mm]},thick] (0,0)--(1.15,0);
  \draw[-{Stealth[length=2mm]},thick] (0,0)--(0,1.15);
  \draw[line width=0.9pt] (1,0) arc (0:90:1);
  \draw[black!70] (0,0)--({cos(45)},{sin(45)})--({cos(45)},0);
  \ptc{45}{$M$}
  \node[font=\scriptsize,below,inner sep=2pt] at ({cos(45)},0) {$H$};
  \node[font=\scriptsize,below left,inner sep=2pt] at (0,0) {$O$};
\end{tikzpicture}
\end{minipage}

\exo[Démontrer les valeurs en $\frac{\pi}{3}$ et $\frac{\pi}{6}$]
\begin{minipage}[c]{0.62\linewidth}\raggedright
$A$ est le point de coordonnées $(1\,;0)$, $M$ l'image de $\frac{\pi}{3}$, $N$ celle de
$\frac{\pi}{6}$ et $N'$ celle de $-\frac{\pi}{6}$.
\begin{enumerate}[label=\alph*)]
  \item Montrer que le triangle $OAM$ est équilatéral.
  \item Sa hauteur issue de $M$ est aussi une médiane : en déduire
        $\cos\left(\frac{\pi}{3}\right)$, puis calculer $\sin\left(\frac{\pi}{3}\right)$.
  \item Montrer que le triangle $ONN'$ est équilatéral.
  \item En déduire la longueur $NN'$, puis $\sin\left(\frac{\pi}{6}\right)$ et
        $\cos\left(\frac{\pi}{6}\right)$.
\end{enumerate}
\end{minipage}\hfill
\begin{minipage}[c]{0.34\linewidth}\centering
\begin{tikzpicture}[x=1.05cm,y=1.05cm]
  \cerclenu
  \draw[black!70] (0,0)--({cos(60)},{sin(60)})--(1,0);
  \draw[black!70] (0,0)--({cos(30)},{sin(30)})--({cos(30)},{-sin(30)})--cycle;
  \ptc{60}{$M$}\ptc{30}{$N$}\ptc{-30}{$N'$}
  \path[coulCourbeB] (1,0) node[croix]{};
  \node[font=\scriptsize,below right,inner sep=1pt] at (1,0) {$A$};
\end{tikzpicture}
\end{minipage}

\partiefiche{Angles associés}

\exo[Lire les symétries]
\begin{minipage}[c]{0.56\linewidth}\raggedright
Le point $M$ est l'image de $\frac{\pi}{3}$ : ses coordonnées sont
$\left(\frac{1}{2}\,;\frac{\sqrt{3}}{2}\right)$. Pour chaque point $N$, $R$ et $S$ :
\begin{enumerate}[label=\alph*)]
  \item dire, par une phrase, de quelle symétrie il est l'image de $M$ ;
  \item donner le nombre réel dont il est l'image, entre $-\pi$ et $\pi$ ;
  \item donner son cosinus et son sinus.
\end{enumerate}
\end{minipage}\hfill
\begin{minipage}[c]{0.42\linewidth}\centering
\begin{tikzpicture}[x=1.15cm,y=1.15cm]
  \draw[lecture] ({cos(60)},{sin(60)})--({cos(60)},{-sin(60)});
  \draw[lecture] ({cos(60)},{sin(60)})--({-cos(60)},{sin(60)});
  \draw[lecture] ({cos(60)},{sin(60)})--({-cos(60)},{-sin(60)});
  \cerclerep
  \ptc{60}{$M$}\ptc{-60}{$N$}\ptc{120}{$R$}\ptc{240}{$S$}
\end{tikzpicture}
\end{minipage}

\exo[À l'aide d'un angle associé]
Déterminer la valeur exacte de chaque nombre, à l'aide d'un angle associé :
a)~$\cos\left(\frac{5\pi}{6}\right)$ ; b)~$\sin\left(\frac{4\pi}{3}\right)$ ;
c)~$\cos\left(\frac{3\pi}{4}\right)$ ; d)~$\sin\left(-\frac{\pi}{6}\right)$ ;
e)~$\cos\left(\frac{7\pi}{4}\right)$ ; f)~$\sin\left(\frac{11\pi}{6}\right)$.

\exo[Un angle hors du tableau]
On admet que $\cos\left(\frac{\pi}{12}\right)=\frac{\sqrt{6}+\sqrt{2}}{4}$ et
$\sin\left(\frac{\pi}{12}\right)=\frac{\sqrt{6}-\sqrt{2}}{4}$.
\begin{enumerate}[label=\alph*)]
  \item Vérifier que $\frac{11\pi}{12}=\pi-\frac{\pi}{12}$ et que
        $\frac{5\pi}{12}=\frac{\pi}{2}-\frac{\pi}{12}$.
  \item En déduire la valeur exacte de $\cos\left(\frac{11\pi}{12}\right)$ et de
        $\cos\left(\frac{5\pi}{12}\right)$.
  \item Calculer de même $\sin\left(\frac{13\pi}{12}\right)$,
        $\sin\left(\frac{7\pi}{12}\right)$ et $\cos\left(-\frac{\pi}{12}\right)$.
\end{enumerate}

\exo[Le dire en phrase]
Les points $M$, $N$ et $P$ du cercle trigonométrique sont les images de $\frac{\pi}{7}$,
de $\frac{8\pi}{7}$ et de $\frac{5\pi}{14}$.
\begin{enumerate}[label=\alph*)]
  \item Vérifier que $\frac{8\pi}{7}=\pi+\frac{\pi}{7}$, puis compléter par une phrase :
        « $M$ et $N$ sont symétriques par rapport à \ldots{} ; leurs cosinus sont
        \ldots{} et leurs sinus sont \ldots{} ».
  \item Écrire $\cos\left(\frac{8\pi}{7}\right)$ à l'aide de $\cos\left(\frac{\pi}{7}\right)$.
  \item Vérifier que $\frac{5\pi}{14}=\frac{\pi}{2}-\frac{\pi}{7}$, puis écrire
        $\cos\left(\frac{5\pi}{14}\right)$ à l'aide de $\sin\left(\frac{\pi}{7}\right)$.
\end{enumerate}

\partiefiche{Simplifier une expression à l'aide des angles associés}

\exo[Simplifier]
Pour tout nombre réel $x$, simplifier :
\begin{enumerate}[label=\alph*)]
  \item $A(x)=\cos(\pi-x)+\cos(-x)$
  \item $B(x)=\sin(\pi+x)+\sin(\pi-x)+\sin(-x)$
  \item $C(x)=\cos\left(\frac{\pi}{2}-x\right)-\sin(\pi-x)+\cos\left(\frac{\pi}{2}+x\right)$
  \item $D(x)=2\cos(x+2\pi)+\cos(\pi+x)$
\end{enumerate}

\exo[Faire apparaître un tour]
Pour tout nombre réel $x$, simplifier, en faisant apparaître un angle associé :
a)~$\sin(3\pi+x)$ ; b)~$\cos(x-\pi)$, qui s'écrit $\cos\big({-}(\pi-x)\big)$ ;
c)~$\cos\left(\frac{5\pi}{2}-x\right)$ ; d)~$\sin\left(x-\frac{\pi}{2}\right)$.

\exo[Toujours le même nombre]
Montrer que, pour tout nombre réel $x$ :
\begin{enumerate}[label=\alph*)]
  \item $\sin^{2}\left(\frac{\pi}{2}-x\right)+\sin^{2}(-x)=1$ ;
  \item $\big(\cos(x)+\sin(x)\big)^{2}+\big(\cos(x)-\sin(x)\big)^{2}=2$ ;
  \item $\cos^{2}(\pi+x)+\cos^{2}\left(\frac{\pi}{2}+x\right)=1$.
\end{enumerate}

\exo[Des produits]
Pour tout nombre réel $x$, simplifier :
\begin{enumerate}[label=\alph*)]
  \item $A(x)=\cos(x+\pi)\sin\left(\frac{\pi}{2}-x\right)-\sin^{2}(-x)$ ;
  \item $B(x)=\sin\left(\frac{\pi}{2}+x\right)\cos(-x)+\cos\left(\frac{\pi}{2}-x\right)\sin(\pi-x)$.
\end{enumerate}

\partiefiche{Le cercle trigonométrique complet}

Le cercle complet sert aux parties VII et VIII.
\begin{center}
\begin{tikzpicture}[x=2.75cm,y=2.75cm]
  \foreach \a in {30,45,60}{%
    \draw[lecture] ({cos(\a)},{-sin(\a)})--({cos(\a)},{sin(\a)});
    \draw[lecture] ({-cos(\a)},{-sin(\a)})--({-cos(\a)},{sin(\a)});
    \draw[lecture] ({-cos(\a)},{sin(\a)})--({cos(\a)},{sin(\a)});
    \draw[lecture] ({-cos(\a)},{-sin(\a)})--({cos(\a)},{-sin(\a)});}
  \cerclebase
  \path[coulCourbeB] (1,0) node[croix]{} (0,1) node[croix]{} (-1,0) node[croix]{} (0,-1) node[croix]{};
  \node[font=\footnotesize,inner sep=1pt] at (1.1,0.09) {$0$};
  \node[font=\footnotesize,inner sep=1pt] at (-1.1,0.09) {$\pi$};
  \node[font=\footnotesize,inner sep=1pt] at (-0.12,1.12) {$\frac{\pi}{2}$};
  \node[font=\footnotesize,inner sep=1pt] at (-0.16,-1.12) {$-\frac{\pi}{2}$};
  \foreach \a/\l in {30/\frac{\pi}{6},45/\frac{\pi}{4},60/\frac{\pi}{3},
     120/\frac{2\pi}{3},135/\frac{3\pi}{4},150/\frac{5\pi}{6},
     -30/-\frac{\pi}{6},-45/-\frac{\pi}{4},-60/-\frac{\pi}{3},
     -120/-\frac{2\pi}{3},-135/-\frac{3\pi}{4},-150/-\frac{5\pi}{6}}{%
    \path[coulCourbeB] ({cos(\a)},{sin(\a)}) node[croix]{};
    \node[font=\footnotesize,inner sep=1pt] at ({1.14*cos(\a)},{1.14*sin(\a)}) {$\l$};}
  \foreach \v/\l in {0.5/\frac{1}{2},0.866/\frac{\sqrt{3}}{2},-0.5/-\frac{1}{2},-0.866/-\frac{\sqrt{3}}{2}}{%
    \draw (\v,0.02)--(\v,-0.02);
    \node[below,font=\scriptsize,inner sep=1pt,outer sep=2pt,coulCourbeB,fill=white] at (\v,0) {$\l$};
    \draw (0.02,\v)--(-0.02,\v);
    \node[left,font=\scriptsize,inner sep=1pt,outer sep=2pt,coulCourbeC,fill=white] at (0,\v) {$\l$};}
  \foreach \v/\l in {0.707/\frac{\sqrt{2}}{2},-0.707/-\frac{\sqrt{2}}{2}}{%
    \draw (\v,0.02)--(\v,-0.02);
    \node[above,font=\scriptsize,inner sep=1pt,outer sep=2pt,coulCourbeB,fill=white] at (\v,0) {$\l$};
    \draw (0.02,\v)--(-0.02,\v);
    \node[right,font=\scriptsize,inner sep=1pt,outer sep=2pt,coulCourbeC,fill=white] at (0,\v) {$\l$};}
\end{tikzpicture}
\end{center}

\exo[Retrouver la mesure principale]
Déterminer la mesure principale de l'angle orienté de mesure $x$ tel que :
\begin{enumerate}[label=\alph*)]
  \item $\cos(x)=\frac{\sqrt{2}}{2}$ et $\sin(x)=\frac{\sqrt{2}}{2}$ ;
  \item $\cos(x)=-\frac{\sqrt{3}}{2}$ et $\sin(x)=-\frac{1}{2}$ ;
  \item $\cos(x)=\frac{1}{2}$ et $\sin(x)=-\frac{\sqrt{3}}{2}$ ;
  \item $\cos(x)=-1$ et $\sin(x)=0$ ;
  \item $\cos(x)=-\frac{\sqrt{2}}{2}$ et $\sin(x)=\frac{\sqrt{2}}{2}$ ;
  \item $\cos(x)=0$ et $\sin(x)=-1$.
\end{enumerate}

\exo[Lire le cercle complet]
Compléter le tableau à l'aide du cercle complet, après avoir cherché la mesure
principale quand il le faut.
\begin{center}
{\renewcommand{\arraystretch}{2.4}\setlength{\tabcolsep}{3pt}
\begin{tabular}{|c|c|c|c|c|c|c|}
\hline
$x$ & $-\dfrac{2\pi}{3}$ & $\dfrac{3\pi}{4}$ & $-\dfrac{\pi}{6}$ & $\dfrac{5\pi}{3}$ & $\dfrac{11\pi}{4}$ & $-\dfrac{3\pi}{2}$\\
\hline
$\cos(x)$ & & & & & & \\
\hline
$\sin(x)$ & & & & & & \\
\hline
\end{tabular}}
\end{center}

\exo[Possible ou non ?]
Existe-t-il un nombre réel $x$ tel que :
\begin{enumerate}[label=\alph*)]
  \item $\cos(x)=\frac{\sqrt{3}}{2}$ et $\sin(x)=-\frac{\sqrt{3}}{2}$ ?
  \item $\cos(x)=-\frac{1}{2}$ et $\sin(x)=-\frac{\sqrt{3}}{2}$ ?
  \item $\cos(x)=\frac{3}{5}$ et $\sin(x)=\frac{4}{5}$ ?
\end{enumerate}
Quand il en existe un sur le cercle complet, donner sa mesure principale.

\exo[Un hexagone régulier]
L'hexagone régulier $ABCDEF$ est inscrit dans le cercle trigonométrique, dans le sens
direct ; $A$ est le point de coordonnées $(1\,;0)$.
\begin{enumerate}[label=\alph*)]
  \item Donner la mesure principale associée à chaque sommet.
  \item Donner les coordonnées exactes de chaque sommet.
\end{enumerate}

\partiefiche{Équations avec des cosinus et des sinus}

\exo[Sur $\intof{-\pi}{\pi}$]
Résoudre sur l'intervalle $\intof{-\pi}{\pi}$, à l'aide du cercle complet :
a)~$\cos(x)=\frac{\sqrt{2}}{2}$ ; b)~$\sin(x)=-\frac{\sqrt{3}}{2}$ ;
c)~$\cos(x)=-\frac{1}{2}$ ; d)~$\sin(x)=0$ ; e)~$\cos(x)=-1$ ;
f)~$\sin(x)=\frac{5}{4}$.

\exo[Sur un autre intervalle]
Résoudre, à l'aide du cercle complet :
\begin{enumerate}[label=\alph*)]
  \item $\sin(x)=-\frac{1}{2}$ sur $\intff{0}{2\pi}$ ;
  \item $\cos(x)=-\frac{\sqrt{3}}{2}$ sur $\intff{0}{2\pi}$ ;
  \item $\cos(x)=0$ sur $\intff{-2\pi}{2\pi}$ ;
  \item $\sin(x)=\frac{\sqrt{2}}{2}$ sur $\intff{-2\pi}{2\pi}$ ;
  \item $\cos(x)=1$ sur $\intff{0}{6\pi}$.
\end{enumerate}

\exo[Se ramener à $\cos(x)=a$]
Résoudre :
\begin{enumerate}[label=\alph*)]
  \item $2\cos(x)-1=0$ sur $\intof{-\pi}{\pi}$ ;
  \item $2\sin(x)+\sqrt{3}=0$ sur $\intff{0}{2\pi}$ ;
  \item $\sqrt{2}\cos(x)+1=0$ sur $\intof{-\pi}{\pi}$ ;
  \item $4\sin^{2}(x)=1$ sur $\intof{-\pi}{\pi}$.
\end{enumerate}

\exo[Point sur un cercle de rayon $4$]
Un point $M$ tourne sur le cercle de centre $O$ et de rayon $4$ ; ses coordonnées sont
$\big(4\cos(x)\,;4\sin(x)\big)$. Pour quelles valeurs de $x$ dans $\intof{-\pi}{\pi}$ :
\begin{enumerate}[label=\alph*)]
  \item l'abscisse de $M$ vaut-elle $2$ ?
  \item l'ordonnée de $M$ vaut-elle $-2\sqrt{2}$ ?
  \item le point $M$ a-t-il pour coordonnées $\big({-}2\sqrt{3}\,;2\big)$ ?
\end{enumerate}

\partiefiche{Fonctions cosinus et sinus}

\exo[Tableaux de variation]
Dresser le tableau de variation, à l'aide des courbes du cours :
\begin{enumerate}[label=\alph*)]
  \item de la fonction cosinus sur $\intff{-2\pi}{0}$ ;
  \item de la fonction sinus sur $\intff{-\frac{5\pi}{2}}{-\frac{\pi}{2}}$ ;
  \item de la fonction cosinus sur $\intff{3\pi}{5\pi}$ ;
  \item de la fonction sinus sur $\intff{\frac{3\pi}{2}}{\frac{7\pi}{2}}$.
\end{enumerate}

\exo[Démontrer une période]
Ces fonctions sont définies sur $\R$. Démontrer que :
\begin{enumerate}[label=\alph*)]
  \item $f(x)=5-2\cos(x)$ est périodique de période $2\pi$ ;
  \item $g(x)=\cos(3x)$ est périodique de période $\frac{2\pi}{3}$ ;
  \item $h(x)=\sin\left(\frac{x}{4}\right)$ est périodique de période $8\pi$ ;
  \item $k(x)=\sin^{2}(x)$ est périodique de période $\pi$ (angles associés).
\end{enumerate}

\exo[La période de $\cos(kx)$]
Soit $k$ un nombre réel strictement positif et $f(x)=\cos(kx)$, pour tout nombre réel
$x$.
\begin{enumerate}[label=\alph*)]
  \item Montrer que $f\left(x+\frac{2\pi}{k}\right)=f(x)$.
  \item En déduire une période de $\cos(5x)$, puis, de la même façon, une période de
        $\sin\left(\frac{x}{3}\right)$.
  \item Sur le secteur, la tension s'écrit $u(t)=325\sin(100\pi t)$, avec $t$ en
        secondes. Quelle est sa période ? Combien de périodes y a-t-il en une seconde ?
\end{enumerate}

\exo[La fonction $|\sin|$]
Soit $g(x)=|\sin(x)|$, pour tout nombre réel $x$.
\begin{enumerate}[label=\alph*)]
  \item À l'aide d'un angle associé, montrer que $g$ est périodique de période $\pi$.
  \item Sur $\intff{0}{\pi}$, $g(x)=\sin(x)$ : dresser le tableau de variation de $g$
        sur $\intff{0}{\pi}$, puis sur $\intff{\pi}{2\pi}$.
\end{enumerate}

\partiefiche{Résolution graphique d'équations et d'inéquations}

Les courbes ci-dessous servent aux exercices 39 à 42 ; un carreau vaut $\frac{\pi}{6}$ en
abscisse et $\frac{1}{2}$ en ordonnée.
\begin{center}
\begin{tikzpicture}[x=0.6cm,y=0.6cm]
  \reperetrigo{-12}{12}\courbetrigo{cos}
  \gradpi{-12/-2\pi,-6/-\pi,6/\pi,12/2\pi}
  \node[coulCourbe,font=\scriptsize,fill=white,inner sep=1pt] at (2.3,1.25) {$y=\cos(x)$};
\end{tikzpicture}

\medskip
\begin{tikzpicture}[x=0.6cm,y=0.6cm]
  \reperetrigo{-12}{12}\courbetrigo{sin}
  \gradpi[above]{-12/-2\pi,-6/-\pi,6/\pi,12/2\pi}
  \node[coulCourbe,font=\scriptsize,fill=white,inner sep=1pt] at (3.4,1.25) {$y=\sin(x)$};
\end{tikzpicture}
\end{center}

\exo[Avec le cosinus]
Résoudre graphiquement sur $\intff{-2\pi}{2\pi}$ :
a)~$\cos(x)=\frac{1}{2}$ ; b)~$\cos(x)=-1$ ; c)~$\cos(x)\ge 0$ ;
d)~$\cos(x)<-\frac{1}{2}$.

\exo[Avec le sinus]
Résoudre graphiquement sur $\intff{-2\pi}{2\pi}$ :
a)~$\sin(x)=-\frac{1}{2}$ ; b)~$\sin(x)=1$ ; c)~$\sin(x)<0$ ;
d)~$\sin(x)\ge\frac{\sqrt{3}}{2}$, la droite d'équation $y=\frac{\sqrt{3}}{2}$ passant
par le point de coordonnées $\left(\frac{\pi}{3}\,;\frac{\sqrt{3}}{2}\right)$.

\exo[Deux conditions à la fois]
Pour quelles valeurs de $x$ dans $\intff{0}{2\pi}$ a-t-on à la fois :
\begin{enumerate}[label=\alph*)]
  \item $\cos(x)\ge 0$ et $\sin(x)\le 0$ ?
  \item $\sin(x)\ge\frac{1}{2}$ et $\cos(x)\ge 0$ ?
  \item $\cos(x)\le-\frac{1}{2}$ et $\sin(x)>0$ ?
\end{enumerate}

\exo[Combien de solutions ?]
Sans les calculer, donner le nombre de solutions sur $\intff{-2\pi}{2\pi}$ de chaque
équation : a)~$\cos(x)=0{,}3$ ; b)~$\sin(x)=-0{,}8$ ; c)~$\cos(x)=1$ ;
d)~$\sin(x)=1{,}2$.

\end{multicols*}

\end{document}
